% SPDX-License-Identifier: GPL-3.0-or-later
% Copyright (C) 2026 Alessandro Greco
%% Testo alternativo per l'accessibilita' (usato solo nell'EPUB).
\altfig{Confronto fra tre generazioni di onion routing: nelle prime due il client prepara un'unica cipolla di cifratura per l'intero percorso, mentre in Tor il circuito si estende un salto alla volta, con uno scambio di chiavi separato per ogni relay.}
% Le tre generazioni dell'onion routing dell'NRL
\begin{tikzpicture}[font=\sffamily\scriptsize,
  nodo/.style={relay,minimum size=0.65cm,font=\sffamily\tiny},
  eti/.style={anchor=east,align=right,font=\sffamily\small}]
  % prima generazione
  \node[eti] at (-0.5,0) {\textbf{1ª gen.}\\(1996)};
  \node[client,minimum width=1.1cm,minimum height=0.6cm,font=\sffamily\tiny,align=center] (c1) at (0.6,0) {client\\+ nodo};
  \foreach \i in {1,...,5} \node[nodo] (a\i) at (0.6+1.05*\i,0) {};
  \draw[cellaf] (c1) -- (a1); \foreach \i/\j in {1/2,2/3,3/4,4/5} \draw[cellaf] (a\i) -- (a\j);
  \draw[decorate,decoration={brace,amplitude=3pt}] (1.3,0.45) -- node[above=3pt,font=\sffamily\tiny] {5 nodi fissi; un'unica cipolla costruisce tutto il circuito} (5.95,0.45);
  % seconda generazione
  \node[eti] at (-0.5,-1.6) {\textbf{2ª gen.}};
  \node[client,minimum width=1.1cm,minimum height=0.6cm,font=\sffamily\tiny] (c2) at (0.6,-1.6) {client};
  \foreach \i in {1,...,4} \node[nodo] (b\i) at (0.6+1.05*\i,-1.6) {};
  \draw[cellaf] (c2) -- (b1); \foreach \i/\j in {1/2,2/3,3/4} \draw[cellaf] (b\i) -- (b\j);
  \node[anchor=west,font=\sffamily\tiny,align=left] at (5.4,-1.6) {client separato dal nodo,\\lunghezza variabile, mixing (abbandonato)};
  % terza generazione
  \node[eti] at (-0.5,-3.2) {\textbf{3ª gen.}\\Tor (2002)};
  \node[client,minimum width=1.1cm,minimum height=0.6cm,font=\sffamily\tiny] (c3) at (0.6,-3.2) {client};
  \foreach \i in {1,2,3} \node[nodo] (d\i) at (0.6+1.05*\i,-3.2) {};
  \draw[cellaf] (c3) -- node[above,font=\sffamily\tiny] {1} (d1);
  \draw[cellaf] (d1) -- node[above,font=\sffamily\tiny] {2} (d2);
  \draw[cellaf] (d2) -- node[above,font=\sffamily\tiny] {3} (d3);
  \node[anchor=west,font=\sffamily\tiny,align=left] at (4.3,-3.2) {3 nodi, costruzione \emph{un salto alla volta}:\\uno scambio di chiavi per salto, forward secrecy};
\end{tikzpicture}
